% !TEX root = hospital_pharmacy_slides_v7_v2.tex
% !TEX program = xelatex
\documentclass{beamer}

\usepackage[UTF8]{ctex}
\usepackage{graphicx}
\usepackage{booktabs}
\usepackage{amsmath}

\usetheme{default}
\usecolortheme{default}
\setbeamertemplate{navigation symbols}{}
\setbeamertemplate{footline}[frame number]

\title{Pharmaceutical Reform and Pharmacy Entry}
\subtitle{County-level Results, v2: One-Year Effective Policy Timing}
\author{Wenxia Chen}
\date{June 26, 2026}

\newcommand{\countyvtwoimage}[1]{%
\IfFileExists{../hospital_with_pharmacy/integrated_results/county_v2/#1}{%
\includegraphics[width=\textwidth]{../hospital_with_pharmacy/integrated_results/county_v2/#1}%
}{%
\fbox{\parbox[c][0.42\textheight][c]{0.95\textwidth}{\centering \scriptsize Missing county v2 image:\\#1}}%
}%
}

\begin{document}

\frame{\titlepage}

\section{County-level Results}

\begin{frame}{Empirical Setting}
\small
\begin{itemize}\setlength\itemsep{0.6em}
    \item Unit of analysis: county-by-year panel.
    \item Controlled sample: 2012--2019, with population, nighttime lights, and road density.
    \item v2 timing treats the reform as effective in the year after the policy year:
    \[
        Reform_{ct}=\mathbf{1}\{t\geq PolicyYear_c+1\}.
    \]
    \item Standard errors are clustered at the county level.
    \item TWFE and CSDID event-study figures use the same outcome definitions and event-time window as the five-estimator robustness figures.
\end{itemize}
\end{frame}

\begin{frame}{TWFE Baseline Results: Count and ln(count)}
\centering
\scriptsize
\setlength{\tabcolsep}{8pt}
\begin{tabular}{lcc}
\toprule
& Count & ln(count) \\
\midrule
Post policy
& 3.739$^{*}$ & 0.031$^{*}$ \\
& (1.565) & (0.014) \\
\midrule
County FE & Yes & Yes \\
Year FE & Yes & Yes \\
Controls & Yes & Yes \\
\midrule
N & 17,463 & 15,940 \\
R$^2$ & 0.842 & 0.876 \\
\bottomrule
\end{tabular}

\vspace{0.8em}
\raggedright
\tiny Controls include log population, log nighttime lights, and log road density.
\end{frame}

\begin{frame}{TWFE Event-study Results}
\centering
\begin{columns}
\begin{column}{0.49\textwidth}
\countyvtwoimage{twfe_event_pharmacy_count.png}
\end{column}
\begin{column}{0.49\textwidth}
\countyvtwoimage{twfe_event_ln_pharmacy.png}
\end{column}
\end{columns}

\vspace{0.3em}
\tiny Relative year $-1$ is omitted. County and year fixed effects are absorbed.
\end{frame}

\begin{frame}{CSDID Results}
\centering
\scriptsize
\setlength{\tabcolsep}{8pt}
\begin{tabular}{lcc}
\toprule
& Count & ln(count) \\
\midrule
ATT
& -2.608 & -0.040 \\
& (1.819) & (0.026) \\
\midrule
Estimator & CSDID & CSDID \\
Controls & Yes & Yes \\
Control group & Not-yet-treated & Not-yet-treated \\
\bottomrule
\end{tabular}

\vspace{0.8em}
\raggedright
\tiny CSDID uses Callaway--Sant'Anna group-time ATT with doubly robust inverse probability weighting and county-clustered standard errors.
\end{frame}

\begin{frame}{CSDID Event-study Results}
\centering
\begin{columns}
\begin{column}{0.49\textwidth}
\countyvtwoimage{csdid_event_pharmacy_count.png}
\end{column}
\begin{column}{0.49\textwidth}
\countyvtwoimage{csdid_event_ln_pharmacy.png}
\end{column}
\end{columns}

\vspace{0.3em}
\tiny Same event window as the TWFE and five-estimator event-study figures; relative year $-1$ is omitted.
\end{frame}

\begin{frame}{Five-estimator Event-study Robustness}
\centering
\begin{columns}
\begin{column}{0.49\textwidth}
\countyvtwoimage{event_5estimators_count.png}
\end{column}
\begin{column}{0.49\textwidth}
\countyvtwoimage{event_5estimators_ln_count.png}
\end{column}
\end{columns}

\vspace{0.3em}
\tiny Estimators: TWFE, Callaway--Sant'Anna CSDID, Sun--Abraham interaction-weighted, Borusyak--Jaravel--Spiess imputation, and stacked event study.
\end{frame}

\begin{frame}{Continuous DID: OLS and IV Baseline Results}
\centering
\scriptsize
\setlength{\tabcolsep}{4.5pt}
\begin{tabular}{lcccc}
\toprule
& \multicolumn{2}{c}{OLS} & \multicolumn{2}{c}{IV: 1959 stock} \\
\cmidrule(lr){2-3}\cmidrule(lr){4-5}
& Count & ln(count) & Count & ln(count) \\
\midrule
Post $\times$ ln(public hospitals)
& 20.308$^{***}$ & -0.028 & 79.687$^{***}$ & 0.107 \\
& (3.460) & (0.021) & (13.734) & (0.083) \\
\midrule
County FE & Yes & Yes & Yes & Yes \\
Year FE & Yes & Yes & Yes & Yes \\
Controls & Yes & Yes & Yes & Yes \\
\midrule
N & 17,463 & 15,940 & 14,241 & 12,886 \\
R$^2$ / K-P F & 0.843 & 0.876 & 100.124 & 147.517 \\
\bottomrule
\end{tabular}

\vspace{0.6em}
\raggedright
\tiny IV instruments interact post policy with the 1959 epidemiology-station stock.
\end{frame}

\begin{frame}{Continuous DID Event-study: OLS}
\centering
\begin{columns}
\begin{column}{0.49\textwidth}
\countyvtwoimage{continuous_did_event_count.png}
\end{column}
\begin{column}{0.49\textwidth}
\countyvtwoimage{continuous_did_event_ln_count.png}
\end{column}
\end{columns}

\vspace{0.3em}
\tiny Event-time indicators are interacted with baseline public-hospital intensity.
\end{frame}

\begin{frame}{Continuous DID Event-study: IV}
\centering
\begin{columns}
\begin{column}{0.49\textwidth}
\countyvtwoimage{continuous_did_event_iv_count.png}
\end{column}
\begin{column}{0.49\textwidth}
\countyvtwoimage{continuous_did_event_iv_ln_count.png}
\end{column}
\end{columns}

\vspace{0.3em}
\tiny Dynamic IV first stages are weak, so these figures are diagnostic rather than the main IV evidence.
\end{frame}

\begin{frame}{Takeaways}
\small
\begin{itemize}\setlength\itemsep{0.65em}
    \item TWFE shows smaller positive count and log-count effects after retaining 2012 policy-year counties.
    \item CSDID static ATT remains negative for count and log count, so staggered-timing robustness does not support the TWFE increase.
    \item Five-estimator event-study figures place TWFE and CSDID in the same event-time comparison set.
    \item Continuous DID OLS remains positive for pharmacy counts; static IV estimates are larger and have strong first stages.
    \item IV event-study plots are retained as diagnostics; the main IV evidence is the static continuous-DID specification.
\end{itemize}
\end{frame}

\end{document}
